\documentclass[a4paper,11pt]{article}

\usepackage[ngerman]{babel}
\usepackage[top=3cm,bottom=3cm,left=3cm,right=3cm]{geometry} %NICHT IN KEMIE2


%\usepackage{microtype} %schöneres typesetting  mit [stretch=10] für nur 1% anstelle von 2%
%\usepackage{lmodern} %Latin Modern FONT
%\usepackage{pifont} %schönere Font



\usepackage{amsmath}
\usepackage{nicefrac}
\usepackage{amssymb}
%\usepackage{mathtools}

\usepackage{titlesec} %Um Section, etc bearbeiten zu können
\usepackage{setspace}
\usepackage{enumitem} %enumerate
\usepackage{multicol}
\usepackage[many]{tcolorbox}
\usepackage{hyperref}
\usepackage{arydshln}

%\renewcommand{\ttdefault}{lmtt} %Schriftart wird geändert zu Latin Modern Typewriter


\hypersetup{hidelinks,
			colorlinks=true,
			linkcolor=black,
			citecolor=miudiutex,
			urlcolor=miugray2
			}

\setlength{\parindent}{0pt}
\setcounter{section}{0}


\titlespacing*{\section}{0pt}{20mm}{6mm}
\titlespacing*{\subsection}{0pt}{6mm}{3mm}
\titlespacing*{\subsubsection}{0pt}{3mm}{0mm}



%\definecolor{miudiutex}{HTML}{2f3d39}
\definecolor{miudiutex}{HTML}{1d2926}
\definecolor{miugray}{HTML}{b4cccf}
\definecolor{miugray2}{HTML}{4d7365}
\definecolor{red}{HTML}{cf1f5c}
\definecolor{green}{HTML}{00A36C}
\definecolor{delphinidin}{HTML}{315794}
\definecolor{pelargonidin}{HTML}{b33807}
\definecolor{cyanidin}{HTML}{ba0746}
\definecolor{petunidin}{HTML}{b56fed}



\raggedcolumns %wichtig für Multicolumns


%================================
% SYMBOLS
%================================

\newcommand{\RA}{$\rightarrow$~}
\newcommand{\RL}{$\rightleftharpoons$~}
\newcommand{\HARP}{\ding{226}~}
%$\xrightarrow{2}$ %Pfeil mit 2 drüber


%================================
% SHORT COMMANDS (BOXES ; ETC)
%================================


\newcommand{\notiz}[1]{\begin{notizz}{}{}\begin{center}
			{\color{miugray2!50!miudiutex}#1} \end{center}\end{notizz}}
\newcommand{\zit}[2]{\begin{zitat*}{#1}{}\small #2 \end{zitat*}}

\newcommand{\frage}[2]{\begin{qst*}{#1}{}\small #2 \end{qst*}}

\newcommand{\unten}{\end{center} \tcblower \begin{center}}


\newcommand*\circled[1]{\tikz[baseline=(char.base)]{
            \node[shape=circle,draw,inner sep=0.5pt,miudiutex] (char) {\tiny #1};}}
            
\newcommand{\miusquare}{{\color{miugray2!50}\scriptsize $\blacksquare$}~}
\newcommand{\miusquarp}{{\color{petunidin!50}\scriptsize $\blacktriangleright$}~}
\newcommand{\niu}{\item[\miusquare]}
\newcommand{\nip}{\item[\miusquarp]}

%%%%% ABSAETZE

\newcommand{\pps}{\par \vspace*{1mm}}
\newcommand{\pp}{\par \vspace*{3mm}}
\newcommand{\ppbig}{\par \vspace*{8mm}}
\newcommand{\pphuge}{\par \vspace*{10mm}}

%================================
% FRAGE BOX
%================================

\tcbuselibrary{theorems,skins,hooks}
\newtcbtheorem{qst}{Frage:}
{%
	enhanced,
	breakable,
	colback = gray!15!white,
	frame hidden,
	boxrule = 0sp,
	borderline west = {2.5pt}{0pt}{red!70},
	sharp corners,
	detach title,
	before upper = \tcbtitle\par\smallskip,
	coltitle = black,
	fonttitle = \bfseries\sffamily,
	description font = \mdseries,
	separator sign none,
	segmentation style={solid, gray!50!black},
}
{th}




%================================
% ZITAT BOX
%================================

\tcbuselibrary{theorems,skins,hooks}
\newtcbtheorem{zitat}{Zitat}
{%
	enhanced,
	breakable,
	colback = gray!15!white,
	frame hidden,
	boxrule = 0sp,
	borderline west = {2.5pt}{0pt}{black!70},
	sharp corners,
	detach title,
%	before upper = \tcbtitle\par\smallskip,
	coltitle = gray!50!black,
	fonttitle = \bfseries\sffamily,
	description font = \mdseries,
%	separator sign none,
	segmentation style={solid, gray!50!black},
}
{th}



%================================
% SIMPLE SCHWARZER RAHMEN BOX
%================================


\newcommand{\boxfr}[1]{
	\begin{tcolorbox}[
	drop shadow southeast,
	enhanced,
	colframe=black!50,
	colback=white,
	boxrule = 1pt, 
%	toprule=0pt, bottomrule=0pt,rightrule=0pt,leftrule=0pt,
	boxsep=5pt,
	arc=1pt,
	]
	\begin{center}
	#1
	\end{center}
	\end{tcolorbox}
}

\definecolor{miudiutex}{HTML}{1d2926}
\definecolor{miugray}{HTML}{b4cccf}
\definecolor{miugray2}{HTML}{4d7365}
\definecolor{white2}{HTML}{f5f7f8}
\definecolor{red}{HTML}{cf1f5c}
\definecolor{green}{HTML}{00A36C}
\definecolor{delphinidin}{HTML}{315794}
\definecolor{uniblau}{HTML}{004e9f}
\definecolor{unigelb}{HTML}{fcba00}
\definecolor{unigrau}{HTML}{909085}
\definecolor{pelargonidin}{HTML}{b33807}
\definecolor{cyanidin}{HTML}{ba0746}
\definecolor{paeonidin}{HTML}{d15ca6}
\definecolor{petunidin}{HTML}{b56fed}
\definecolor{malvidin}{HTML}{8a2d8c}

\newcommand{\metermilli}{\mskip3mu \text{mm}}
\newcommand{\cm}{\mskip3mu \text{cm}}
\newcommand{\m}{ \text{m}}
\newcommand{\lite}{ \text{l}}
\newcommand{\kg}{ \text{kg}}
\newcommand{\g}{ \text{g}}
\newcommand{\km}{ \text{km}}
\newcommand{\h}{ \text{h}}
\newcommand{\s}{ \text{s}}
\newcommand{\tonne}{ \text{t}}
\usepackage[utf8]{inputenc}
\usepackage[T1]{fontenc}
\usepackage{lmodern}

\usepackage[
    activate={true,nocompatibility},
    final,
    tracking=true,
    kerning=true,
    spacing=true,
    factor=1100,
    stretch=30,
    shrink=20
]{microtype}
\usepackage{pdfpages}
\usepackage{awesomebox}
\usepackage{marginnote}
\tcbuselibrary{documentation}
\usepackage{sidenotes}
\usepackage{import}
\usepackage{xifthen}
\usepackage{transparent}
\usepackage[table]{xcolor}


\newcommand{\abbicon}{%
  \protect\tikz[baseline=0.2mm,scale=0.8]{%
    % Das blaue Quadrat mit abgerundeten Ecken
    \fill[uniblau, rounded corners=1.2pt] (0,0) rectangle (0.8em, 0.8em);
    % Der weiße Kreis in der Mitte
    \fill[white] (0.4em, 0.4em) circle (0.12em);
  }%
} 
\usepackage[labelfont={bf},font={color=miudiutex,small},figurename={\abbicon $~$Abbildung}]{caption}


\newcommand{\bckgrnd}{
\AddToHookNext{shipout/background}{
\begin{tikzpicture}[remember picture, overlay]
 \draw[draw=none,fill=uniblau!70, shift={(current page.north west)}] (-3,0) rectangle (23,-3.75);
 \draw[draw=none,fill=miugray, shift={(current page.north east)}] (-7,0) rectangle (3,-3.75);
\end{tikzpicture}
}
}



\newcommand{\incfig}[2]{%
\def\svgwidth{#2\columnwidth}
\import{C://LaTeX-Files/pngs/}{#1.pdf_tex}
}

\renewcommand{\textbf}[1]{\textsf{{\bfseries {#1}}}}



\titleformat{\section}[hang]
		{\LARGE \color{uniblau!75!black} \sffamily \bfseries}
		{\thesection}
		{6mm}
		{}
		[]
\titleformat{\subsection}[hang]
		{\large \sffamily \bfseries \color{black}}
		{\thesubsection}
		{4mm}
		{{\color{miugray}$\blacksquare ~$}}

\titleformat{\subsubsection}[block]
		{\normalsize \bfseries \sffamily \color{miudiutex}}%
		{\normalsize \color{black} \thesubsubsection}
		{2mm}
		{{\color{miugray}$\blacksquare ~$}}
		[ \color{black}]
		
\titleformat{\paragraph}
		{\normalsize \bfseries \sffamily \color{black}}
		{\footnotesize \theparagraph}
		{1mm}
		{}

\pagestyle{empty}
\titlespacing*{\section}{0pt}{20mm}{12mm}
\titlespacing{\section}{0pt}{20mm}{12mm}

\titlespacing*{\subsection}{0pt}{14mm}{4mm}
\titlespacing{\subsection}{0pt}{14mm}{4mm}

\titlespacing*{\subsubsection}{0pt}{6mm}{3mm}
\titlespacing{\subsubsection}{0pt}{6mm}{3mm}

\titlespacing*{\paragraph}{0pt}{6mm}{2mm}
\titlespacing{\paragraph}{0pt}{6mm}{2mm}


\let\oldmarginfigure\marginfigure
\let\endoldmarginfigure\endmarginfigure

\let\oldmarginfigure\marginfigure
\let\endoldmarginfigure\endmarginfigure

\renewenvironment{marginfigure}[1][0pt] % [1] steht für 1 Argument, [0pt] ist der Standardwert
  {\oldmarginfigure[#1]\captionsetup{labelfont={sf,bf,color=uniblau!75!black},font={color=miudiutex,small}}}
  {\endoldmarginfigure}
  
  
  
  
\renewcommand{\labelitemi}{\color{uniblau!70}$\bullet$}
\renewcommand{\labelitemii}{\color{uniblau!70}$\cdot$}
\newcommand{\plmtsquare}{{\color{uniblau!70}$\blacksquare~$}}
\newcommand{\splmt}[1]{{\color{uniblau!75!black} \sffamily #1}}
\newcommand{\fplmt}[1]{{\color{uniblau!75!black} \sffamily \bfseries #1}}

\newcommand{\antwort}[1]{
\begin{tcolorbox}[
	enhanced,
	enforce breakable,
	standard jigsaw,
	boxrule=0.7pt,
	colframe=black,
	sharp corners,
	boxsep=5pt,
	title=\bfseries \sffamily Antwort,
	opacityback=0,
	coltitle=miudiutex,
	attach boxed title to top text left={yshift=-\tcboxedtitleheight/2},
	boxed title style={colback=white,colframe=white},
	bottomrule at break=0pt,
	toprule at break=0pt,
	pad at break=16mm,
]
	\begin{center}
	#1
	\end{center}
	\end{tcolorbox}
}




\rowcolors{2}{miugray!50}{miugray!50}
\arrayrulecolor{white} % Weiße Gitterlinien
\setlength{\arrayrulewidth}{1.5pt} % Etwas dickere Linien
\renewcommand{\arraystretch}{1.4} % Mehr Zeilenhöhe für bessere Lesbarkeit}


%\setlist[itemize]{itemsep=0mm}
\setlist[enumerate]{itemsep=0mm}

\usetikzlibrary{patterns, arrows.meta, calc}
\usepackage{pgfplots}
\pgfplotsset{compat=1.18}

\newif\ifshowme
%\showmetrue




\begin{document}
\newgeometry{top=1.8cm,bottom=4cm,left=1.5cm,right=7.5cm,marginparsep=-2.00cm,marginparwidth=4.75cm}


\bckgrnd


\begin{center}
\LARGE
\color{white}
\textbf{Übung 06} \par \vspace*{2mm}
\large \color{miugray} \textbf{Prozessbezogene Lebensmitteltechnologie}
\end{center}
\par 
\vspace*{6mm}


\color{miudiutex}

\section*{Haltbarmachung I}

\subsection*{Aufgabe 1} Beim Erhitzen einer Sporensuspension bei 112 °C werden folgende Sporenzahlen nach unterschiedlichen Erhitzungszeiten festgestellt:


\pp
\begin{center}\begin{tabular}{r|r}
\rowcolor{miugray}
Zeit [min] & Sporenzahl \\ \hline
0 &1.000.000 \\
4 &110.000 \\
8 &12.000 \\
12 &1.200 \\
\end{tabular}\end{center}


\ppbig
\marginnote[]{ \small Der \fplmt{D-Wert} beschreibt die Zeit in Minuten, die benötigt wird, um die Keimzahl um eine log-Stufe bzw. um eine Zehnerpotenz zu reduzieren.}

Bestimmen Sie den \splmt{D-Wert} sowohl 
\begin{enumerate}[label=\alph*)]
\item mittels semi-logarithmischer graphischer Auftragung (nutzen Sie dafür den vorgegebenen Graphen) 
\item rechnerisch
\end{enumerate}

\begin{center}
\resizebox{\textwidth}{!}{
\begin{tikzpicture}
    \begin{semilogyaxis}[
        width=12cm,
        height=8cm,
        xmin=0, xmax=15,       % Grenzen der X-Achse
        ymin=1000, ymax=1500000,    % Grenzen der Y-Achse (3 Dekaden: 10^0 bis 10^3)
        enlarge y limits={upper, value=0.05},
        axis lines=left,      % Achsenlinien nur links und unten
		ymajorgrids=true,     % NUR horizontale Haupt-Gitterlinien zeichnen
        yminorgrids=true,     % NUR horizontale Zwischen-Gitterlinien zeichnen
        major grid style={black!60, thin}, % Stil der Haupt-Gitterlinien (1, 10, 100...)
        minor grid style={black!30, thin}, % Stil der Zwischen-Gitterlinien (2, 3, 4...)
        minor y tick num=8,   % 8 Zwischenstriche pro Dekade
        xtick={0,1,2,3,4,5,6,7,8,9,10,11,12,13,14},
%        xticklabels={\empty}, % Versteckt die X-Achsen-Zahlen (wie im Bild)
%        yticklabels={\empty}, % Versteckt die Y-Achsen-Zahlen (wie im Bild)
        tick align=inside,   % Achsenstriche zeigen nach außen
        ylabel=Keimzahl,
        xlabel=t,
%        axis background/.style={fill=miugray!50}
    ]
%\addplot[color=blue, very thick, domain=0:7, samples=100] {1.5^x};
    % --- Hier kannst du deinen Graphen einfügen ---
    % Beispiel:
\ifshowme

     \addplot[color=uniblau, thick, mark=*] coordinates {(0,1000000) (4, 110000) (8, 12000) (12, 1200)};

\else
\fi

    
    \end{semilogyaxis}
\end{tikzpicture}
}
\end{center}

\ifshowme
\antwort{
\begin{align*}
D = \frac{4 ~\text{min}}{\log (\frac{10^6}{1,1 \cdot 10^5})} = 4,17 ~\text{min}
\end{align*}
}
\else
\fi



\subsection*{Aufgabe 2} Für eine Sporensuspension wurden folgende D-Werte bestimmt:

\begin{center}\begin{tabular}{r|r}
\rowcolor{miugray}
Temperatur [°C] & D-Wert [min] \\ \hline
104 &27,5 \\
107 &14,5 \\
110 &7,5 \\
113 &4,0 \\
116 &2,2 \\

\end{tabular}\end{center}
\pp

\marginnote[]{\small Der \fplmt{z-Wert} gibt an, wie stark die temperaturabhängige Änderung des D-Werts ist. Ein hoher z-Wert bedeutet, dass die Änderung durch die Temperatur sehr stark ist.}
Bestimmen Sie den \splmt{z-Wert}
\begin{enumerate}[label=\alph*)]
\item mittels semi-logarithmischer graphischer Auftragung (nutzen Sie dafür den vorgegebenen Graphen) 
\item rechnerisch
\end{enumerate}

\begin{center}
\resizebox{\textwidth}{!}{
\begin{tikzpicture}
    \begin{semilogyaxis}[
        width=12cm,
        height=7cm,
        xmin=100, xmax=120,       % Grenzen der X-Achse
        ymin=1, ymax=150,    % Grenzen der Y-Achse (3 Dekaden: 10^0 bis 10^3)
        enlarge y limits={upper, value=0.05},
        axis lines=left,      % Achsenlinien nur links und unten
		ymajorgrids=true,     % NUR horizontale Haupt-Gitterlinien zeichnen
        yminorgrids=true,     % NUR horizontale Zwischen-Gitterlinien zeichnen
        major grid style={black!60, thin}, % Stil der Haupt-Gitterlinien (1, 10, 100...)
        minor grid style={black!30, thin}, % Stil der Zwischen-Gitterlinien (2, 3, 4...)
        minor y tick num=8,   % 8 Zwischenstriche pro Dekade
        xtick={102,104,106,108,110,112,114,116,118},
%        xticklabels={\empty}, % Versteckt die X-Achsen-Zahlen (wie im Bild)
%        yticklabels={\empty}, % Versteckt die Y-Achsen-Zahlen (wie im Bild)
        tick align=inside,   % Achsenstriche zeigen nach außen
        ylabel=D-Wert (in min),
        xlabel=Temperatur (°C)
    ]
    
\ifshowme

     \addplot[color=uniblau, thick, mark=*] coordinates {(104,27.5) (107,14.5) (110,7.5) (113, 4) (116,2.2)};
     
\else
\fi


    
    \end{semilogyaxis}
\end{tikzpicture}
}
\end{center}

\ifshowme
\antwort{
\begin{align*}
z &= \frac{T_2 - T_1}{\log(D_1) - \log (D_2)} \\
z &= \frac{107 - 104}{\log(27,5) - \log (14,5)} \\
 &= 10,79
\end{align*}

}
\else
\fi



\subsection*{Aufgabe 3}
In einem Lebensmittel liegt die Anfangskeimzahl bei 1.000 Sporen pro Gramm. Es wird ein 12D-Konzept angewandt. Das Lebensmittel wird in 500 g-Packungen abgefüllt. 
\pp
In wie vielen Packungen ist nach der Erhitzung theoretisch noch eine Spore enthalten?

\ifshowme
\antwort{

Sporen pro Gramm:
\begin{align*}
1000 ~(\text{Sporen} /\text{g}) \cdot 10^{-12} = 10^{-9} ~(\text{Sporen} /\text{g})
\end{align*}
Sporen pro 500 g - Packung:
\begin{align*}
10^{-9} \cdot 500 ~\text{g} = \frac{1}{2.000.000}
\end{align*}
In einer von 2 Millionen Packungen ist noch eine Spore enthalten.
}
\else
\fi


\subsection*{Aufgabe 4}
Die Anfangskeimbelastung liegt bei 10.000 Sporen pro Konserve. Es wird 50 min lang bei
113°C erhitzt. Der D-Wert liegt bei dieser Temperatur bei 4 min. 
\pp 
Wie hoch ist das Risiko, dass eine Packung eine Spore enthält?


\ifshowme
\antwort{
Formel aus Food Processing Handbook, S.35:
\begin{align}
\frac{\text{Heating Time}}{D_T} &= \log(\frac{N_0}{N}) 
\end{align}
Zunächst potenzieren wir um den Logarithmus umstellen zu können:
\begin{align*}
10^{\frac{\text{Heating Time}}{D_T}} &= \frac{N_0}{N} \\
N &= \frac{N_0}{10^{\frac{\text{Heating Time}}{D_T}}}
\end{align*}
Werte einsetzen:
\begin{align*}
N &= \frac{10.000}{10^{\frac{\text{50 ~\text{min}}}{4 ~\text{min}}}} \\
 &= 3,16 \cdot 10^{-9}
\end{align*}
Anzahl der Konserven, die geöffnet werden müssen um eine Spore zu finden:
\begin{align*}
\frac{1}{3,16 \cdot 10^{-9}} = 3,16 \cdot 10^{8}
\end{align*}
}
\else
\fi




















 

\subsection*{Aufgabe 5}
Wie lange muss erhitzt werden, um bei T = 75 °C einen Endkeimgehalt bzgl.
\textit{Streptococcus spp.} von 10$^{-3}$ pro Behältnis zu erreichen?
\par 
\boxfr{
Konserveninhalt: 450 g \par 
Anfangskeimgehalt: 10$^2$/g \par 
D$_T^{Z}$ = D$_{80 ~\text{°C}}^{8,5 ~\text{°C}}$ = 0,5 min \par 
}

\pp

\ifshowme
\antwort{
D-Wert der neuen Temperatur ausrechnen:
\begin{align*}
D_1 &= D_2 \cdot 10^{\frac{T_2-T_1}{z}} \\
D_1 &= 0,5 ~\text{min}  \cdot 10^{\frac{80-75}{8,5}} \\
	&= 1,94
\end{align*}


Anfangskeimgehalt einer 450 g-Konserve: 
\begin{align*}
10^2 \cdot 450 = 45.000 
\end{align*} 
Zeit die nötig ist, um auf 10$^{-3}$ Sporen pro Behältnis zu reduzieren:
\begin{align*}
\text{Heating Time} &= D_T \cdot \log(\frac{N_0}{N}) \\
\text{Heating Time} &= 1,94 \cdot \log(\frac{45.000}{10^{-3}}) \\
 &= 14,85 ~ \text{min}
\end{align*}

}
\else
\fi



\subsection*{Aufgabe 6}
In der folgenden Tabelle fehlen einige Angaben. Bitte füllen Sie diese Lücken aus. Nehmen Sie an, dass es sich hierbei nur um eine Reaktion handelt, sodass die Spalten
zueinander in einer Beziehung stehen. \par 

\begin{center}
\begin{spacing}{1.2}
\ifshowme
\antwort{
\begin{tabular}{|l|l|l|}
D-Wert & Temperatur & Zeit \\
\hline
D$_{80\text{°C}}^{5 \text{°C}}$ = 0,5 min & T = 80 °C & t = 0,5 min \\
\hline
D$_{70\text{°C}}^{5 \text{°C}}$ = 50 min & T = 70°C  & t = 50 min \\
\hline
D$_{85\text{°C}}^{5 \text{°C}}$ = 0,05 min & T = 85 °C & t = 0,05 min \\
\hline
D$_{100\text{°C}}^{5 \text{°C}}$ = 0,00005 min & T = 100°C &t = 0,00005 min \\
\hline
\end{tabular}
}
\else
\begin{tabular}{|l|l|l|}
\rowcolor{miugray}
D-Wert & Temperatur & Zeit \\
\hline
D$_{80\text{°C}}^{5 \text{°C}}$ = 0,5 min & T = 80 °C & t = 0,5 min \\
\hline
D$_{??}^{??}$ = 50 min & T = ??  & t = ?? min \\
\hline
D$_{??}^{??}$ = ?? min & T = 85 °C & t = ?? min \\
\hline
D$_{??}^{??}$ = ?? min & T = ?? &t = 0,00005 min \\
\hline
\end{tabular}
\fi
\end{spacing}
\end{center}

\subsection*{Aufgabe 7}
Sie haben für Ihren Erhitzungsprozess einen D-Wert von D$_{121,1}$ = 3 min gegeben. Wie sieht
der D-Wert bei einer Temperatur von 112 °C aus, wenn ein z-Wert von 10 °C gegeben ist?

\ifshowme
\antwort{
\begin{align*}
D_1 &= D_2 \cdot 10^{\frac{T_2-T_1}{z}} \\
D_1 &= 3 \cdot 10^{\frac{121,1-112}{10}} \\
 &= 24,38 ~\text{min}
\end{align*}
}
\else
\fi



\subsection*{Aufgabe 8}
Für die Abtötung von Mikroorganismen in einem Fleischprodukt soll der z-Wert ermittelt
werden. Es liegen folgende Daten vor:
\pp
\begin{center}
\fbox{
\parbox[c][1cm][c]{0.4\textwidth}{
T1 = 50 °C ; T2 = 65 °C \par D1 = 180 s ; D2 = 5 s
}
}
\end{center}

\ifshowme
\antwort{
\begin{align*}
D_1 &= D_2 \cdot 10^{\frac{T_2-T_1}{z}} \\
\log(\frac{D_1}{D_2}) &= \frac{T_2-T_1}{z} \\
z &= \frac{T_2 - T_1}{\log(\frac{D_1}{D_2})} \\
z &= \frac{65 - 50}{\log(\frac{180}{5})} \\ 
 &= 9,64  ~\text{°C}
\end{align*}
}
\else
\fi



\subsection*{Aufgabe 9}
Ein Lebensmittel wird für 2 Minuten schlagartig auf 150°C erhitzt. Berechnen Sie den F$_0$-Wert, wenn der z-Wert 15 K beträgt.
\ifshowme
\antwort{
\begin{align*}
F_1 &= F_2 \cdot 10^{\frac{T_2-T_1}{z}} \\
F_1 &= 2 ~\text{min} \cdot 10^{\frac{150-121,1}{15}} \\
 &= 168,926 ~\text{min}
\end{align*}
}
\else
\fi


\subsection*{Aufgabe 10}
Wie groß ist der Letalitätswert einer einminütigen, schlagartigen Erhitzung bei 128 °C
bzw. 118 °C und einem z-Wert von 6 K?

\ifshowme
\antwort{
\begin{align*}
L &= 10^{\frac{T_2-121,1}{z}} \\ \\
L_{128} &= 10^{\frac{128-121,1}{6}} = 14,125 \\
L_{118} &= 10^{\frac{118-121,1}{6}} = 0,304
\end{align*}
}
\else
\fi


 

\subsection*{Aufgabe 11}
Sie erhalten für den Erhitzungsprozess ihres Produktes für die Qualitätssicherung die
Vorgabe, dass der F$_0$-Wert maximal 5 Minuten betragen darf. Der z-Wert liegt bei 10 K.
Um zu prüfen, ob die Anforderung erfüllt wird, haben Sie in der folgenden Tabelle die
Daten Ihres Erhitzungsprozesses gesammelt. Wird die Anforderung der
Qualitätssicherung erfüllt?

\begin{center}\begin{tabular}{r|r}
\rowcolor{miugray}
Temperatur [°C] & Zeit [s] \\ \hline
95 & 10 \\
110 & 15 \\
130 & 42 \\
120 & 35 \\
90 & 10 \\

\end{tabular}\end{center}

\ifshowme
\antwort{
\begin{align*}
F_0  &= &F_t \cdot 10^{\frac{T_2-T_1}{z}} \qquad && \\ \\
F_01 &= &10 ~\text{s} \cdot 10^{\frac{95-121,1}{10}}  \qquad &= 0,025 ~\text{s} \\
F_02 &= &15 ~\text{s} \cdot 10^{\frac{110-121,1}{10}} \qquad &= 1,16 ~\text{s} \\
F_03 &= &42 ~\text{s} \cdot 10^{\frac{130-121,1}{10}} \qquad &= 326,02 ~\text{s} \\
F_04 &= &35 ~\text{s} \cdot 10^{\frac{120-121,1}{10}} \qquad &= 27,17 ~\text{s} \\
F_05 &= &10 ~\text{s} \cdot 10^{\frac{90-121,1}{10}}  \qquad &= 0,008 ~\text{s} \\
\end{align*}

Summe aller Zeiten:
\begin{align*}
\sum_{i=1}^{5} F_{0i} = 5,91 ~\text{min}
\end{align*}
}
\else
\fi







\end{document}
